%%% POST-RECORD v1 BEGIN
% format: hexagon-post-format/1
% record_id: H01 (field_experts/HEXAGON_CANDIDATES.md); return HODGE_SIGN_final_20261121T0057Z, findings HS-01..HS-03, SG-08..SG-11
% version: 1
% date: 2026-11-21
% title: "Deligne's purity theorem without monodromy, from an effective trace-formula remainder, Parseval and tensor powers"
% kind: new proof
% claim_status: proved
% statement_plain: >
%   Let X be a separated scheme of finite type over a finite field F_q and F a constructible l-adic sheaf on X whose
%   local Frobenius eigenvalues at every closed point x have absolute value at most |k(x)|^(w/2). Then every Frobenius
%   eigenvalue on the compactly supported cohomology H^i_c(X, F) has absolute value at most q^((w+i)/2). If X is smooth
%   and proper and F is lisse and pointwise pure of weight w, then H^i(X, F) is pure of weight w+i.
% extends: "gives a new proof of Deligne, Weil I, Theoreme 1.6, and of the absolute-value content of Weil II, Theoreme 3.3.1, over a point"
% known_before: "Deligne 1974 (Weil I) and 1980 (Weil II); other proofs by Laumon 1987 (Fourier transform) and Kedlaya 2006 (p-adic)"
% new_here: "a proof that uses no monodromy, no Lefschetz pencil, no Rankin-Selberg step and no theory of weights; its analytic input is Parseval for a bilinear character sum"
% main_statements: [Theorem A = Theorem 6.1, Corollary B = Corollary 6.2]
% hard_step: "Proposition 4.1 with the choice of hypersurface in Theorem 5.1"
% external_results: Katz, JAMS 6 (1993), Corollary 3.4.1 (generic weak Lefschetz) and Corollary 3.4.3 / Lemma 3.5 (generic smoothness); standard etale cohomology (trace formula, Kunneth, affine vanishing, Poincare duality, localization)
% produced_by:
%   - {role: question, model: "Claude Opus 5.5", effort: "max", date: 2026-11-20}
%   - {role: proof, model: "GPT-6 Astra", effort: "ChatGPT Pro mode", date: 2026-11-21}
%   - {role: text of this post, model: "Claude Opus 5.5", effort: "max", date: 2026-11-21}
% checks: see the Verification record
% human_understanding: none
% text_origin: primarily AI-generated
% formalization: none
% licence: CC0-1.0
% repository: https://github.com/KokunoYumeto/grothendieck-montpellier-archive/tree/main/research/grh_20261120/deligne_sign_only
% hexagon:
%   subject: math.AG
%   cross_lists: [math.NT]
%   msc: [14F20, 14G15, 11G25, 11T23]
%   tags: [Weil conjectures, Riemann hypothesis over finite fields, etale cohomology, weak Lefschetz, character sums, tensor power trick]
%%% POST-RECORD v1 END
\documentclass[11pt,a4paper]{article}
\usepackage[T1]{fontenc}
\usepackage{lmodern}
\usepackage{microtype}
\usepackage[margin=2.6cm]{geometry}
\usepackage{amsmath,amssymb,amsthm}
\usepackage{booktabs,tabularx,longtable,array}
\usepackage{enumitem}
\usepackage{xcolor}
\PassOptionsToPackage{hyphens}{url}
\usepackage[hidelinks]{hyperref}
\emergencystretch=2em

\theoremstyle{plain}
\newtheorem{maintheorem}{Theorem}
\renewcommand{\themaintheorem}{\Alph{maintheorem}}
\newtheorem{maincorollary}[maintheorem]{Corollary}
\newtheorem{theorem}{Theorem}[section]
\newtheorem{lemma}[theorem]{Lemma}
\newtheorem{proposition}[theorem]{Proposition}
\newtheorem{corollary}[theorem]{Corollary}
\theoremstyle{definition}
\newtheorem{cited}[theorem]{Cited result (not proved here)}
\newtheorem{definition}[theorem]{Definition}
\theoremstyle{remark}
\newtheorem{remark}[theorem]{Remark}

\newcommand{\uses}[1]{\par\noindent{\small\textsf{Uses:} #1}\par\smallskip}
\newcommand{\step}[2]{\par\smallskip\noindent\textbf{Step #1} (#2).\ \ignorespaces}
\newcolumntype{L}[1]{>{\raggedright\arraybackslash}p{#1}}
\newcolumntype{Y}{>{\raggedright\arraybackslash}X}
\newcommand{\fixedhead}[1]{\section*{#1}}

\newcommand{\Fq}{\mathbf F_q}
\newcommand{\Fqb}{\overline{\mathbf F}_q}
\newcommand{\Ql}{\overline{\mathbf Q}_\ell}
\newcommand{\A}{\mathbf A}
\newcommand{\F}{\mathcal F}
\newcommand{\Tr}{\operatorname{Tr}}
\newcommand{\Frob}{\operatorname{Frob}}

\title{Deligne's purity theorem without monodromy, from an effective trace-formula remainder, Parseval and tensor powers}
\author{Entered at deposit}
\date{21 November 2026 \quad Version 1 \quad CC0 1.0}

\begin{document}
\maketitle

\begin{abstract}
\noindent
Let $X$ be a separated scheme of finite type over a finite field with $q$ elements, and $F$ a constructible
$\ell$-adic sheaf on $X$ whose Frobenius eigenvalues at each closed point $x$ have absolute value at most
$|\kappa(x)|^{w/2}$. Then every Frobenius eigenvalue on the compactly supported cohomology $H^i_c(X,F)$ has absolute
value at most $q^{(w+i)/2}$; if $X$ is smooth and proper and $F$ is lisse and pointwise pure of weight $w$, then
$H^i(X,F)$ is pure of weight $w+i$. These are theorems of Deligne (La conjecture de Weil I, 1974, and II, 1980),
reproved by Laumon (1987) and by Kedlaya (2006). What is new is the proof: it uses no monodromy, no Lefschetz pencil, no
Rankin-Selberg step and no theory of weights. On a power $U^k$ of a smooth affine $U$, cut by a generic hypersurface
built from a bilinear form in Noether coordinates, the generic weak Lefschetz
theorem (Deligne's proof in the appendix of Katz 1993) makes the remainder of the point count the trace of an actual
Frobenius module containing the $k$-th tensor power of the middle cohomology of $U$; Parseval bounds this remainder at
square-root size up to a factor independent of $k$, and letting $k$ grow finishes. The proof was found by GPT-6 Astra (ChatGPT
Pro mode), working out a remark of L. Alp\"oge (2026). It was checked step by step by Claude Opus 5.5 at effort high and
re-proved independently by Claude Fable 5.2 at effort max, whose one correction (the source of the Lefschetz input) is
adopted; the text is by Claude Opus 5.5 at effort max. No human has checked it.
\end{abstract}

% ==================================================================================
\section{Result}
% ==================================================================================

\subsection{Statement}

Fix a prime power $q$, a prime $\ell\nmid q$ and a field embedding $\iota:\Ql\hookrightarrow\mathbf C$; absolute values
of $\ell$-adic numbers are taken through $\iota$. For a scheme $X$ over $\Fq$, $H^i_c(X,\F)$ means $H^i_c(X_{\Fqb},\F)$
with the action of the geometric Frobenius $F$. A constructible $\Ql$-sheaf $\F$ on $X$ has \emph{pointwise weight
$\le w$} ($w\in\mathbf R$) if for every closed point $x$ every eigenvalue $\lambda$ of the geometric Frobenius $\Frob_x$
on $\F_{\bar x}$ satisfies $|\iota\lambda|\le|\kappa(x)|^{w/2}$; it is \emph{pointwise pure of weight $w$} if equality
holds for all of them.

\begin{maintheorem}[$=$ Theorem~\ref{thm:all}]\label{thm:A}
Let $X$ be a separated scheme of finite type over $\Fq$, $w\in\mathbf R$, and $\F$ a constructible $\Ql$-sheaf on $X$ of
pointwise weight $\le w$. Then for every $i$ every eigenvalue $\alpha$ of $F$ on $H^i_c(X,\F)$ satisfies
$|\iota\alpha|\le q^{(w+i)/2}$.
\end{maintheorem}

\begin{maincorollary}[$=$ Corollary~\ref{cor:pure}]\label{cor:B}
Let $X$ be a smooth proper scheme over $\Fq$ and $\F$ a lisse $\Ql$-sheaf on $X$, pointwise pure of weight $w$. Then for
every $i$ every eigenvalue $\alpha$ of $F$ on $H^i(X,\F)$ satisfies $|\iota\alpha|=q^{(w+i)/2}$. In particular, for $X$
smooth and projective, the eigenvalues of $F$ on $H^i(X,\Ql)$ have absolute value $q^{i/2}$ for every embedding $\iota$.
\end{maincorollary}

\noindent\textbf{Instance.} For an elliptic curve $E/\Fq$ and $\F=\Ql$, Corollary~\ref{cor:B} gives $|\iota\alpha|=q^{1/2}$
for the two eigenvalues on $H^1(E,\Ql)$, hence Hasse's bound $|\#E(\Fq)-q-1|\le2\sqrt q$.

\smallskip
\noindent\textbf{Not claimed.} The rationality, functional equation and integrality parts of the Weil conjectures, which
are separate known results; semisimplicity of Frobenius; the hard Lefschetz theorem; the theory of weights for complexes
and its further theorems in Weil II; algebraicity of the eigenvalues of an arbitrary constructible sheaf.

\subsection{Status}

\noindent
\begin{tabularx}{\textwidth}{@{}L{3.2cm}Y@{}}
\toprule
Kind & New proof of known theorems (Deligne, Weil I, Th\'eor\`eme 1.6; the absolute-value content of Weil II, Th\'eor\`eme
3.3.1, over a point).\\
Claim & Proved here, from the cited results of Section~\ref{sec:inputs}.\\
Checks & Three AI checks, all on 21 Nov 2026 (details in the Verification record):
\begin{itemize}[leftmargin=*,nosep]
\item Claude Opus 5.5 at effort high: a targeted proof check of every step.
\item Claude Fable 5.2 at effort max: an independent re-proof of every step, without seeing that check. It made one
correction, the source of the Lefschetz input; Lemma~\ref{lem:good} adopts it.
\item The main session (Claude Opus 5.5 at effort max) rederived Lemma~\ref{lem:parseval},
Proposition~\ref{prop:remainder} and the exponents of Theorem~\ref{thm:middle}.
\end{itemize}\\
Human understanding & None. No human has read or checked the proof.\\
Formalization & None.\\
Open points & None in the statements. The foundational references (SGA 4, SGA $4\frac12$, the elementary part of BBD
\S4) were not reopened; see Limits.\\
\bottomrule
\end{tabularx}

\subsection{Known and new}

\noindent
\begin{tabularx}{\textwidth}{@{}L{4.6cm}Y@{}}
\toprule
Deligne \cite{WeilI}, Th\'eor\`eme 1.6 & Corollary~\ref{cor:B} for $\F=\Ql$; proof by Lefschetz pencils, the monodromy theorem
of Kazhdan--Margulis and Rankin's squaring.\\
Deligne \cite{WeilII}, Th\'eor\`eme 3.3.1 & Theorem~\ref{thm:A} (in the language of mixed sheaves, over any base); proof by
monodromy and the case of curves.\\
Laumon \cite{Laumon}; Kedlaya \cite{Kedlaya} & Other proofs of Weil II: by the $\ell$-adic Fourier transform; by
$p$-adic methods.\\
Scholl \cite{Scholl} & Reduces Weil I to smooth hypersurfaces without Lefschetz pencils, and says an elementary,
Stepanov-style proof of the point-count estimate would be ``extremely interesting''.\\
Katz \cite{KatzACT}, App.~A.5 (proof by Deligne); Cor.~3.4.1 & The generic weak Lefschetz theorem used here, in the form
of Lemma~\ref{lem:good}.\\
Alp\"oge \cite{Alpoge}, footnote 3 & The idea: a sign-only proof of Weil and Deligne from $qN=M^k+R$.\\
\midrule
New here & A complete proof along that idea: the effective remainder (Proposition~\ref{prop:remainder}), the choice of
hypersurface that makes the remainder a bilinear character sum (Theorem~\ref{thm:middle}), and the deduction of
Theorem~\ref{thm:A} and Corollary~\ref{cor:B}.\\
Literature search & On 21 Nov 2026 we ran two web searches, one for new proofs of the Weil conjectures and one for
elementary approaches in higher dimension. They found the alternatives above and no proof along this route. The
independent checker, from what it knows, knew no earlier proof in arbitrary dimension that avoids both a monodromy
theorem and Rankin's method. The search does not settle priority.\\
\bottomrule
\end{tabularx}

\subsection{Proof map}

\textbf{Idea.} Count points on $U^k$ cut by a hypersurface whose equation is a bilinear ``path form'' in Noether
coordinates. The point count is then a bilinear character sum, which Parseval bounds at square-root size whatever the
weights coming from $U$ are. Generic weak Lefschetz makes the difference between this count and its main term the trace
of an actual Frobenius module that contains $H^d_c(U,\F)^{\otimes k}$, so no eigenvalue can cancel against another.
The square-root bound then holds for every eigenvalue, with a loss that does not grow with $k$, and tensor powers remove
the loss.

\textbf{Plan.} Lemma~\ref{lem:power} (effective power sums) $\to$ Lemma~\ref{lem:tensor} (tensor amplification);
Lemma~\ref{lem:parseval} (Parseval along a path); Proposition~\ref{prop:remainder} (the remainder is effective)
$\to$ Theorem~\ref{thm:middle} (the middle degree of a smooth affine scheme) $\to$ Theorem~\ref{thm:all}
($=$ Theorem~\ref{thm:A}) $\to$ Corollary~\ref{cor:pure} ($=$ Corollary~\ref{cor:B}).

\textbf{The hard step.} Proposition~\ref{prop:remainder} together with the choice of hypersurface in
Theorem~\ref{thm:middle}: the hypersurface must satisfy weak Lefschetz (so that the remainder is effective) and have an
equation for which the point count is a bilinear character sum (so that Parseval applies). Generic hyperplane sections
of the graph of the fixed function $f$ (Lemma~\ref{lem:good}) give both at once.

\textbf{External results.} Section~\ref{sec:inputs}: the trace formula, K\"unneth, affine vanishing, Poincar\'e duality
and localization; Noether normalization; the generic weak Lefschetz theorem of Katz's appendix. None involves weights.

\medskip
\noindent
\begin{tabularx}{\textwidth}{@{}L{3.6cm}Y@{}}
\toprule
Statement & Uses\\
\midrule
Lemma~\ref{lem:power} & complex analysis only\\
Lemma~\ref{lem:tensor} & Lemma~\ref{lem:power}\\
Lemma~\ref{lem:parseval} & character orthogonality, Cauchy--Schwarz\\
Lemma~\ref{lem:good} & \ref{cit:katz}; a dimension count\\
Proposition~\ref{prop:remainder} & \ref{cit:trace}, \ref{cit:kunneth}, \ref{cit:artin}, \ref{cit:duality}\\
Theorem~\ref{thm:middle} & Lemmas~\ref{lem:tensor}, \ref{lem:parseval}, \ref{lem:good}, Proposition~\ref{prop:remainder},
\ref{cit:noether}\\
Theorem~\ref{thm:all} & Theorem~\ref{thm:middle}, Lemma~\ref{lem:good}, \ref{cit:artin}, \ref{cit:duality},
\ref{cit:local}\\
Corollary~\ref{cor:pure} & Theorem~\ref{thm:all}, \ref{cit:duality}\\
\bottomrule
\end{tabularx}

% ==================================================================================
\section{Conventions and cited results}\label{sec:inputs}
% ==================================================================================

A \emph{Frobenius module} is a finite-dimensional $\Ql$-vector space with an automorphism. If $V$ is one, with
automorphism $\Phi$, then $\Tr(\Phi^n,V)=\sum_j m_j\mu_j^n$, where the $\mu_j$ are the distinct eigenvalues of $\Phi$ and
$m_j\ge1$ their algebraic multiplicities. For $\F$ on $X/\Fq$ and $x\in X(\mathbf F_{q^n})$, $\Tr(F_{x,n},\F_{\bar x})$
denotes the trace of the geometric Frobenius of $\mathbf F_{q^n}$ at $x$. The following are used; none of them involves
weights.

\begin{cited}[Trace formula; \cite{SGA4h}, Rapport, Th.~3.2; \cite{WeilI}, (1.5.4)]\label{cit:trace}
For $X/\Fq$ separated of finite type, $\F$ constructible and $n\ge1$:
$\sum_{x\in X(\mathbf F_{q^n})}\Tr(F_{x,n},\F_{\bar x})=\sum_i(-1)^i\Tr(F^n,H^i_c(X,\F))$.
\end{cited}

\begin{cited}[K\"unneth; \cite{SGA4}, XVII, 5.4.3]\label{cit:kunneth}
For $X,Y$ separated of finite type over $\Fq$ and constructible $\F,\mathcal G$, there is a Frobenius-equivariant
isomorphism of graded spaces $H^*_c(X\times Y,\F\boxtimes\mathcal G)\cong H^*_c(X,\F)\otimes H^*_c(Y,\mathcal G)$.
\end{cited}

\begin{cited}[Affine vanishing; \cite{SGA4}, XIV, 3.1]\label{cit:artin}
If $X$ is affine of dimension $d$ over an algebraically closed field and $\mathcal G$ is constructible, then
$H^a(X,\mathcal G)=0$ for $a>d$.
\end{cited}

\begin{cited}[Poincar\'e duality; \cite{SGA4}, XVIII, \S3]\label{cit:duality}
If $X/\Fq$ is smooth of pure dimension $d$ and $\F$ is lisse, the cup product
$H^i_c(X,\F)\otimes H^{2d-i}(X,\F^\vee)\to H^{2d}_c(X,\Ql)\cong\Ql(-d)$ is a perfect, Frobenius-equivariant pairing.
Consequently, if $Y\subset X$ is a smooth closed subscheme of pure dimension $d-1$, the dual of the restriction
$H^a(X,\F^\vee)\to H^a(Y,\F^\vee|_Y)$ is a Frobenius-equivariant (Gysin) map
$H^{2d-2-a}_c(Y,\F|_Y)(-1)\to H^{2d-a}_c(X,\F)$. With~\ref{cit:artin}: if $X$ is moreover affine, $H^i_c(X,\F)=0$ for
$i<d$.
\end{cited}

\begin{cited}[Localization; \cite{SGA4}, XVII, \S5]\label{cit:local}
For $U\subset X$ open with closed complement $Z$ and $\F$ constructible, the sequence
$H^i_c(U,\F)\to H^i_c(X,\F)\to H^i_c(Z,\F)$ is exact and Frobenius-equivariant. It is part of the long exact sequence
obtained from $0\to j_!j^*\F\to\F\to i_*i^*\F\to0$ on $X$ ($j:U\to X$, $i:Z\to X$), since
$H^*_c(X,j_!\mathcal G)=H^*_c(U,\mathcal G)$ and $H^*_c(X,i_*\mathcal G)=H^*_c(Z,\mathcal G)$.
\end{cited}

\begin{cited}[Noether normalization; \cite{Stacks}, Tag 00OY]\label{cit:noether}
If $U$ is affine of pure dimension $d$ over a field $k$, there is a finite morphism $\pi:U\to\A^d_k$.
\end{cited}

\begin{cited}[Generic weak Lefschetz; \cite{KatzACT}, Appendix A.5 with A.3--A.4, pp.~219--221]\label{cit:katz}
Let $k$ be an algebraically closed field of characteristic $p\ne\ell$, $X$ a separated $k$-scheme of finite type,
$\varphi:X\to\mathbf P^N$ a quasi-finite morphism and $\mathcal K$ a perverse sheaf on $X$. There is a dense open set of
hyperplanes $H\subset\mathbf P^N$ such that the restriction $H^a(X,\mathcal K)\to H^a(\varphi^{-1}H,\mathcal K)$ is an
isomorphism for $a<-1$ and injective for $a=-1$.
\end{cited}

Katz attributes this proof to Deligne (\cite{KatzACT}, p.~152). It uses the K\"unneth formula with biduality in the
universal situation, generic base change over the space of hyperplanes, proper base change, Verdier duality, the
behaviour of perversity under quasi-finite direct images, and affine vanishing; no weights enter. If $X$ is smooth of
pure dimension $D$ and $\mathcal G$ is lisse, then $\mathcal G[D]$ is perverse, and the conclusion reads: the restriction
$H^a(X,\mathcal G)\to H^a(\varphi^{-1}H,\mathcal G)$ is an isomorphism for $a<D-1$ and injective for $a=D-1$. Katz's
Corollary 3.4.1 (\cite{KatzACT}, p.~181) states the form used below, with a fixed function $f$ and a varying linear part,
but its printed proof runs through his Main Theorem 1.5 and the Fourier transform; Lemma~\ref{lem:good} derives that form
from A.5 instead.

\begin{lemma}[Good hypersurfaces]\label{lem:good}
Let $V$ be a smooth affine scheme over $\Fqb$ of pure dimension $D\ge1$, $y_1,\dots,y_M$ functions on $V$ such that
$\Gamma=(y_1,\dots,y_M):V\to\A^M$ is a closed immersion, and $\mathcal G$ a lisse $\Ql$-sheaf on $V$. Let $\Lambda$ be
either the space of all $(\lambda_1,\dots,\lambda_M,\mu)$ with $\lambda\ne0$, or the space of those with $\lambda_M=1$.
There is a dense open subset $\Omega\subset\Lambda$ such that for $(\lambda,\mu)\in\Omega$ the closed subscheme
$Y=V\cap(\sum_m\lambda_my_m+\mu=0)$ is smooth and, on each connected component of $V$, of pure dimension $D-1$ or empty,
and the restriction $H^a(V,\mathcal G)\to H^a(Y,\mathcal G)$ is an isomorphism for $a<D-1$ and injective for $a=D-1$.
\end{lemma}
\begin{proof}
View $\A^M\subset\mathbf P^M$ as the complement of the hyperplane at infinity. The hyperplanes $H\ne H_\infty$ of
$\mathbf P^M$ are the zero sets of $\sum_m\lambda_my_m+\mu$ with $\lambda\ne0$, up to scaling, and $\Gamma^{-1}H=Y$.

\emph{Lefschetz.} Apply~\ref{cit:katz} to $\varphi=\Gamma$ (a closed immersion followed by an open one, hence quasi-finite)
and $\mathcal K=\mathcal G[D]$: the good hyperplanes form a dense open set $G_1$ of the dual projective space.

\emph{Smoothness, by a dimension count.} Let $I\subset V\times(\mathbf P^M)^\vee$ be the closed set of pairs $(v,H)$ with
$\Gamma(v)\in H$ and $H$ containing the embedded tangent space $d\Gamma(T_vV)$. Since $\Gamma$ is an immersion, this tangent
space has dimension $D$, so the hyperplanes through it form a projective space of dimension $M-D-1$, and
$\dim I\le D+(M-D-1)=M-1$. The closure of the image of $I$ is therefore a proper closed subset of $(\mathbf P^M)^\vee$, and
its complement $G_2$ is dense and open. If $H\in G_2$, the equation of $H$ restricts to a nonzero linear form on the
tangent space at every point of $Y$, so $Y$ is smooth, and on each connected component of $V$ it is of pure dimension
$D-1$ or empty. This count holds in every characteristic, because $\Gamma$ is an immersion (it is the classical Bertini
theorem for immersions; compare \cite{Jouanolou}, I.6.10).

\emph{Charts.} $G_1\cap G_2$ is dense and open. It meets the open chart where $\lambda\ne0$, respectively where
$\lambda_M\ne0$, in a dense open subset, and scaling to $\lambda_M=1$ identifies the latter with a dense open subset of
$\Lambda$.
\end{proof}
\uses{\ref{cit:katz}.}

% ==================================================================================
\section{Three lemmas}
% ==================================================================================

\begin{lemma}[Effective power sums]\label{lem:power}
Let $\lambda_1,\dots,\lambda_s$ be distinct nonzero complex numbers, $c_1,\dots,c_s$ nonzero complex numbers, and
$S_n=\sum_jc_j\lambda_j^n$. If $\limsup_n|S_n|^{1/n}\le B$, then $|\lambda_j|\le B$ for every $j$. This applies when
$|S_n|\le C\,B^n$ for all $n\ge1$.
\end{lemma}
\begin{proof}
The series $G(z)=\sum_{n\ge1}S_nz^n$ converges on $|z|<1/B$. Near $z=0$ it equals the rational function
$R(z)=\sum_jc_j\lambda_jz/(1-\lambda_jz)$. Near $z_j=1/\lambda_j$ the $j$-th term is $-c_jz/(z-z_j)$, which has a simple
pole with residue $-c_j/\lambda_j\ne0$, and the other terms are regular there because the $\lambda_j$ are distinct. So the
poles of $R$ are exactly the points $1/\lambda_j$. By the identity theorem $G=R$ on the disc $|z|<1/B$, where $G$ is
holomorphic; hence no $1/\lambda_j$ lies in that disc, that is, $|\lambda_j|\le B$.
\end{proof}
\uses{nothing beyond complex analysis.}

For an actual Frobenius module the coefficients are algebraic multiplicities $m_j\ge1$, so they are never zero; for a
virtual module (a difference of two modules) an eigenvalue can cancel and escape the bound. This is why
Proposition~\ref{prop:remainder} is needed.

\begin{lemma}[Tensor amplification]\label{lem:tensor}
Let $V$ be a Frobenius module with automorphism $F$, and $a\in\mathbf R$. Let $\mathcal K$ be an infinite set of integers
$k\ge2$, and suppose that for each $k\in\mathcal K$ there are an integer $e_k\ge1$, a Frobenius module $(K_k,\Phi_k)$, a
sign $\varepsilon_k\in\{\pm1\}$, constants $C_k>0$ and $c_k\ge0$ with $c_k/k\to0$ as $k\to\infty$ in $\mathcal K$, and
complex numbers $R_{k,n}$ such that for all $n\ge1$
\[
\varepsilon_kR_{k,n}=\iota\Tr\bigl(F^{e_kn},V^{\otimes k}\bigr)+\iota\Tr\bigl(\Phi_k^n,K_k\bigr),\qquad
|R_{k,n}|\le C_k\,q^{e_kn(ka+c_k)}.
\]
Then every eigenvalue $\alpha$ of $F$ on $V$ satisfies $|\iota\alpha|\le q^a$.
\end{lemma}
\begin{proof}
Fix $k\in\mathcal K$. The right-hand side is the $n$-th power sum of the eigenvalues of the automorphism
$F^{e_k}\oplus\Phi_k$ of $V^{\otimes k}\oplus K_k$, with their algebraic multiplicities. If $v\ne0$ satisfies
$Fv=\alpha v$, then $v^{\otimes k}$ is an eigenvector of $F^{e_k}$ on $V^{\otimes k}$ with eigenvalue $\alpha^{ke_k}$, so
$\alpha^{ke_k}$ occurs among these eigenvalues. Lemma~\ref{lem:power}, applied with $B=q^{e_k(ka+c_k)}$, gives
$|\iota\alpha|^{ke_k}\le q^{e_k(ka+c_k)}$, that is, $|\iota\alpha|\le q^{a+c_k/k}$. Let $k\to\infty$ in $\mathcal K$.
\end{proof}
\uses{Lemma~\ref{lem:power}.}

\begin{lemma}[Parseval along a path]\label{lem:parseval}
Let $E$ be a finite field with $Q$ elements, $\psi:E\to\mathbf C^\times$ a nontrivial additive character, $d\ge1$,
$A=E^d$ and $k\ge2$. Let $b_1,\dots,b_{k-1}:A\times A\to E$ be nondegenerate bilinear forms, $c\in E$, $T\ge0$, and
$t_1,\dots,t_k:A\to\mathbf C$ functions with $|t_i(x)|\le T$ for all $i$ and $x$. Then for every $h\in E^\times$
\[
\Bigl|\sum_{x\in A^k}\ \prod_{i=1}^kt_i(x_i)\ \psi\Bigl(h\Bigl(\sum_{i=1}^{k-1}b_i(x_i,x_{i+1})+c\Bigr)\Bigr)\Bigr|
\le T^k\,Q^{d(k+1)/2}.
\]
\end{lemma}
\begin{proof}
For $i<k$ let $K_i$ be the $A\times A$ matrix $K_i(x,y)=\psi(h\,b_i(x,y))$. For $x,x'\in A$,
\[
(K_iK_i^*)(x,x')=\sum_{y\in A}\psi\bigl(h\,b_i(x,y)-h\,b_i(x',y)\bigr)=\sum_{y\in A}\psi\bigl(h\,b_i(x-x',y)\bigr).
\]
If $x\ne x'$, then $y\mapsto h\,b_i(x-x',y)$ is a nonzero linear form, because $b_i$ is nondegenerate and $h\ne0$, so the
sum vanishes; if $x=x'$ it is $Q^d$. Thus $K_iK_i^*=Q^dI$, and $K_i$ has operator norm $Q^{d/2}$ on $\ell^2(A)$ with the
counting measure. Let $D_i$ be the diagonal matrix with entries $t_i(x)$, of norm at most $T$, and view $t_1$ and $t_k$
as vectors, of norms at most $TQ^{d/2}$. Since $|\psi(hc)|=1$, the sum has the absolute value of
$t_1^{\mathsf T}K_1D_2K_2\cdots D_{k-1}K_{k-1}t_k$. By the Cauchy--Schwarz inequality this is at most
$\|t_1\|\cdot\|K_1\|\cdots\|K_{k-1}\|\cdot\|D_2\|\cdots\|D_{k-1}\|\cdot\|t_k\|
\le(TQ^{d/2})^2\,Q^{d(k-1)/2}\,T^{k-2}=T^kQ^{d(k+1)/2}$.
\end{proof}
\uses{orthogonality of additive characters; the Cauchy--Schwarz inequality.}

% ==================================================================================
\section{The remainder is effective}
% ==================================================================================

\begin{proposition}[Effective remainder]\label{prop:remainder}
Let $r$ be a power of $q$, $U$ a smooth affine scheme over $\mathbf F_r$ of pure dimension $d\ge1$, $\F$ a lisse
$\Ql$-sheaf on $U$, $k\ge1$, $X=U^k$, $D=kd$ and $\F_k=\F^{\boxtimes k}$ on $X$. Let $i:Y\hookrightarrow X$ be a smooth
affine closed subscheme of pure dimension $D-1$ (possibly empty), defined over $\mathbf F_r$, such that over
$\overline{\mathbf F}_r$ the restriction $H^a(X,\F_k^\vee)\to H^a(Y,i^*\F_k^\vee)$ is an isomorphism for $a<D-1$ and
injective for $a=D-1$. For $n\ge1$ let
\[
M_n=\sum_{u\in U(\mathbf F_{r^n})}\Tr(F_{u,n},\F_{\bar u}),\qquad
N_n=\sum_{y\in Y(\mathbf F_{r^n})}\Tr(F_{y,n},(\F_k)_{\bar y}),\qquad R_n=r^nN_n-M_n^k.
\]
Let $F$ be the geometric Frobenius of $\mathbf F_r$. Then
\[
(-1)^{D-1}R_n=\Tr\bigl(F^n,H^d_c(U,\F)^{\otimes k}\bigr)+\Tr(F^n,K),
\]
where
\[
K=\ker\bigl(H^{D-1}_c(Y,i^*\F_k)(-1)\to H^{D+1}_c(X,\F_k)\bigr)
\]
and the map is the Gysin map of~\ref{cit:duality}. In particular $(-1)^{D-1}R_n$ is the $n$-th power sum of the
eigenvalues of $F$ on the actual Frobenius module $H^d_c(U,\F)^{\otimes k}\oplus K$, counted with algebraic multiplicity.
\end{proposition}
\begin{proof}
\step{1}{vanishing} By~\ref{cit:duality}, $H^j_c(U,\F)=0$ for $j<d$, $H^j_c(X,\F_k)=0$ for $j<D$ and
$H^j_c(Y,i^*\F_k)=0$ for $j<D-1$, since $U$, $X$ and $Y$ are smooth and affine of pure dimensions $d$, $D$, $D-1$ and the
sheaves are lisse.

\step{2}{Gysin} By~\ref{cit:duality} applied on $X$ and on $Y$, the dual of the restriction in degree $a$ is the Gysin map
$H^{2D-2-a}_c(Y,i^*\F_k)(-1)\to H^{2D-a}_c(X,\F_k)$. Writing $j=2D-2-a$, the hypothesis says that
$G_j:H^j_c(Y,i^*\F_k)(-1)\to H^{j+2}_c(X,\F_k)$ is an isomorphism for $j>D-1$ and surjective for $j=D-1$. The $G_j$ are
Frobenius-equivariant.

\step{3}{traces} Put $W_a=H^{D-1+a}_c(Y,i^*\F_k)$ and $T_b=H^{D+b}_c(X,\F_k)$ for $a,b\ge0$. The stalk of $\F_k$ at
$y=(y_1,\dots,y_k)$ is $\bigotimes_i\F_{\bar y_i}$, so the trace sum of $\F_k$ over $X(\mathbf F_{r^n})=U(\mathbf F_{r^n})^k$
is $M_n^k$. The Tate twist $(-1)$ multiplies the eigenvalues of $F$ by $r$. By~\ref{cit:trace} and Step~1,
\[
r^nN_n=\sum_{a\ge0}(-1)^{D-1+a}\Tr(F^n,W_a(-1)),\qquad M_n^k=\sum_{b\ge0}(-1)^{D+b}\Tr(F^n,T_b),
\]
hence
\[
(-1)^{D-1}R_n=\sum_{a\ge0}(-1)^a\Tr(F^n,W_a(-1))+\sum_{b\ge0}(-1)^b\Tr(F^n,T_b).
\]

\step{4}{cancellation} For $a\ge1$, $G_{D-1+a}$ identifies $W_a(-1)$ with $T_{a+1}$ (Step~2), and the two terms
$(-1)^a\Tr(F^n,W_a(-1))$ and $(-1)^{a+1}\Tr(F^n,T_{a+1})$ cancel. What remains is
$\Tr(F^n,W_0(-1))-\Tr(F^n,T_1)+\Tr(F^n,T_0)$. Since $G_{D-1}:W_0(-1)\to T_1$ is surjective with kernel $K$, the exact
sequence $0\to K\to W_0(-1)\to T_1\to0$ of Frobenius modules gives $\Tr(F^n,W_0(-1))-\Tr(F^n,T_1)=\Tr(F^n,K)$.

\step{5}{K\"unneth} By~\ref{cit:kunneth}, $T_0=H^D_c(X,\F_k)=\bigoplus_{j_1+\dots+j_k=D}\bigotimes_iH^{j_i}_c(U,\F)$. By
Step~1 only $j_1=\dots=j_k=d$ contributes, so $T_0=H^d_c(U,\F)^{\otimes k}$ as Frobenius modules.
\end{proof}
\uses{\ref{cit:trace}, \ref{cit:kunneth}, \ref{cit:artin}, \ref{cit:duality}.}

No concentration of the cohomology of $U$ or $Y$ in one degree is assumed: the higher degrees cancel through actual
isomorphisms, and the border surjection leaves an actual kernel.

% ==================================================================================
\section{The middle degree of a smooth affine scheme}
% ==================================================================================

\begin{theorem}[Middle degree]\label{thm:middle}
Let $U$ be a smooth affine scheme over $\Fq$ of pure dimension $d\ge1$, $\F$ a lisse $\Ql$-sheaf on $U$, $w\in\mathbf R$
and $C>0$ such that $|\iota\Tr(F_{u,n},\F_{\bar u})|\le C\,q^{nw/2}$ for all $n\ge1$ and $u\in U(\mathbf F_{q^n})$. Then
every eigenvalue $\alpha$ of $F$ on $H^d_c(U,\F)$ satisfies $|\iota\alpha|\le q^{(w+d)/2}$.
\end{theorem}
\begin{proof}
\step{1}{coordinates} By~\ref{cit:noether} choose a finite morphism $\pi:U\to\A^d$ over $\Fq$, and $B$ elements generating
$\mathcal O(U)$ as a module over $\mathcal O(\A^d)$. For every field $L\supset\Fq$ and $x\in\A^d(L)$ the fibre
$\pi^{-1}(x)$ is the spectrum of an $L$-algebra of dimension at most $B$, so it has at most $B$ points with values in $L$.
Choose generators $z_1,\dots,z_s$ of the $\Fq$-algebra $\mathcal O(U)$, so that $z=(z_1,\dots,z_s):U\to\A^s$ is a closed
immersion. Fix $k\ge2$ and a nondegenerate bilinear form $b$ on $\A^d$ over $\Fq$ (for instance the standard dot product).

\step{2}{the family} On $X=U^k$ let $f(u_1,\dots,u_k)=\sum_{i=1}^{k-1}b(\pi(u_i),\pi(u_{i+1}))$, and let
$Z:X\to\A^{sk}$, $u\mapsto(z(u_1),\dots,z(u_k))$, a closed immersion. For $(a,c)\in\A^{sk}\times\A^1$, with
$a=(a_{ij})$, let $Y_{a,c}\subset X$ be the closed subscheme
\[
f(u)+\sum_{i=1}^k\sum_{j=1}^sa_{ij}z_j(u_i)+c=0 .
\]
Over $\Fqb$, $V=X_{\Fqb}$ is smooth and affine of pure dimension $D=kd\ge2$. The graph map
$\Gamma=(Z,f):V\to\A^{sk+1}$ is a closed immersion, because $Z$ is. Apply Lemma~\ref{lem:good} to $\Gamma$, to the lisse
sheaf $\F_k^\vee$, and to the space $\Lambda$ of linear forms whose coefficient on the last coordinate, $f$, is $1$. The
hypersurfaces $\Gamma^{-1}H$ for $H$ in that chart are exactly the $Y_{a,c}$. There is a dense open set $\Omega$ of
parameters $(a,c)$ for which $Y_{a,c}$ is smooth, of pure dimension $D-1$ or empty on each component, and the restriction
$H^a(V,\F_k^\vee)\to H^a(Y_{a,c},\F_k^\vee)$ is an isomorphism for $a<D-1$ and injective for $a=D-1$. The dense open set
$\Omega$ contains a point $(a,c)$ with coordinates in some finite field $\mathbf F_{q^e}$; fix it, put $r=q^e$ and
$Y=Y_{a,c}$. Then $Y$ is defined over $\mathbf F_r$ and satisfies the hypotheses of Proposition~\ref{prop:remainder} over
$\mathbf F_r$. The exponent $e$ and the point $(a,c)$ may depend on $k$.

\step{3}{the remainder as a character sum} Let $n\ge1$, $Q=r^n$, $E=\mathbf F_Q$, and $\psi$ a nontrivial additive
character of $E$. For $u\in X(E)$ write $P(u)$ for the left-hand side of the equation of $Y$, and $\tau(u)=\iota\Tr(F_{u,en},\F_{\bar u})$
for $u\in U(E)$. Since $\sum_{h\in E}\psi(hP)$ is $Q$ if $P=0$ and $0$ otherwise,
\[
Q\,\iota N_n=\sum_{h\in E}\ \sum_{u\in U(E)^k}\ \prod_{i=1}^k\tau(u_i)\ \psi(hP(u)).
\]
Group the $u_i$ by $x_i=\pi(u_i)\in E^d$ and put
$t_{i,h}(x)=\sum_{u\in U(E),\ \pi(u)=x}\tau(u)\,\psi\bigl(h\sum_ja_{ij}z_j(u)\bigr)$. Then the inner sum for a given $h$ is
\[
\psi(hc)\sum_{x\in(E^d)^k}\ \prod_{i=1}^kt_{i,h}(x_i)\ \psi\Bigl(h\sum_{i=1}^{k-1}b(x_i,x_{i+1})\Bigr).
\]
For $h=0$ it equals $\bigl(\sum_{u\in U(E)}\tau(u)\bigr)^k=(\iota M_n)^k$, with $M_n$ formed over $\mathbf F_r$ as in
Proposition~\ref{prop:remainder}. Each fibre of $\pi$ has at most $B$ points in $E$ and $|\tau(u)|\le CQ^{w/2}$, so
$|t_{i,h}(x)|\le BC\,Q^{w/2}$. Lemma~\ref{lem:parseval}, with $T=BC\,Q^{w/2}$ and $b_i=b$, bounds each of the $Q-1$ terms
with $h\ne0$, so
\[
|\iota R_n|=\bigl|Q\,\iota N_n-(\iota M_n)^k\bigr|\le(Q-1)(BC)^kQ^{kw/2}Q^{d(k+1)/2}\le(BC)^k\,q^{en\{k(w+d)/2+d/2+1\}}.
\]

\step{4}{conclusion} By Proposition~\ref{prop:remainder} over $\mathbf F_r$, whose geometric Frobenius is $F^e$,
$(-1)^{D-1}\iota R_n=\iota\Tr(F^{en},H^d_c(U,\F)^{\otimes k})+\iota\Tr((F^e)^n,K)$. So the hypotheses of
Lemma~\ref{lem:tensor} hold for $V=H^d_c(U,\F)$, $a=(w+d)/2$, every $k\ge2$, $e_k=e$, $K_k=K$, $\varepsilon_k=(-1)^{D-1}$,
$C_k=(BC)^k$ and $c_k=d/2+1$. Lemma~\ref{lem:tensor} gives $|\iota\alpha|\le q^{(w+d)/2}$.
\end{proof}
\uses{Lemmas~\ref{lem:tensor}, \ref{lem:parseval} and~\ref{lem:good}; Proposition~\ref{prop:remainder};
\ref{cit:noether}.}

Neither a uniform bound on $e$ nor a bound on the Betti numbers of $Y$ is needed: the limit in $n$ is taken first, for
each fixed $k$.

% ==================================================================================
\section{All degrees, and purity}
% ==================================================================================

\begin{theorem}[$=$ Theorem~\ref{thm:A}]\label{thm:all}
Let $X$ be a separated scheme of finite type over $\Fq$ and $\F$ a constructible $\Ql$-sheaf on $X$ of pointwise weight
$\le w$. Then every eigenvalue $\alpha$ of $F$ on $H^i_c(X,\F)$ satisfies $|\iota\alpha|\le q^{(w+i)/2}$.
\end{theorem}
\begin{proof}
A bound proved after replacing $\Fq$ by $\mathbf F_{q^e}$ descends: the eigenvalues of $F^e$ are the $e$-th powers of
those of $F$, and $|\iota\alpha^e|\le q^{e(w+i)/2}$ gives $|\iota\alpha|\le q^{(w+i)/2}$. Finite base extensions are used
freely below.

\step{1}{smooth affine, lisse} Let $X$ be smooth and affine of pure dimension $d$ and $\F$ lisse of pointwise weight
$\le w$. We prove by induction on $d$ the statement: for every finite extension $\mathbf F_{q^e}$ of $\Fq$, every such
$X$ and $\F$ over $\mathbf F_{q^e}$, and every $i$, the eigenvalues of the geometric Frobenius of $\mathbf F_{q^e}$ on
$H^i_c(X,\F)$ have absolute value at most $q^{e(w+i)/2}$. We write the case $e=1$; the general case is the same with
$q^e$ in place of $q$. If $d=0$, then $H^0_c(X,\F)=\bigoplus_x\F_{\bar x}$ summed over the closed points, where a point of degree $m$
contributes a module induced from $\F_{\bar x}$, on which $F^m$ acts with the eigenvalues of $\Frob_x$; so every
eigenvalue $\alpha$ of $F$ satisfies $|\iota\alpha|^m\le q^{mw/2}$. Let $d\ge1$. For $i<d$, $H^i_c(X,\F)=0$
by~\ref{cit:duality}. For $i=d$ use Theorem~\ref{thm:middle} with $C=\operatorname{rank}\F$. For $i\ge d+1$ choose
closed-immersion coordinates $z:X\to\A^s$. By Lemma~\ref{lem:good} (for $\Gamma=z$, the sheaf $\F^\vee$ and all linear
forms; this is~\ref{cit:katz} with $f=0$), there is, over some finite extension $\mathbf F_{q^{e'}}$, a hyperplane section
$Y=X\cap(\sum_ja_jz_j+c=0)$, smooth and affine of pure dimension $d-1$ or empty, for which the restriction from $X$ to $Y$
on $H^a(-,\F^\vee)$ is an isomorphism for $a<d-1$ and injective for $a=d-1$. Dually~(\ref{cit:duality}), the Gysin map
$H^{i-2}_c(Y,\F|_Y)(-1)\to H^i_c(X,\F)$ is Frobenius-equivariant, surjective for $i=d+1$ and an isomorphism for
$i>d+1$. By the induction hypothesis for $Y$ over $\mathbf F_{q^{e'}}$, the eigenvalues of $F^{e'}$ on
$H^{i-2}_c(Y,\F|_Y)$ have absolute value at most $q^{e'(w+i-2)/2}$; those on its twist $(-1)$, and hence those on
$H^i_c(X,\F)$, are at most $q^{e'}q^{e'(w+i-2)/2}=q^{e'(w+i)/2}$, and the bound for $F$ follows by taking $e'$-th roots.

\step{2}{general $X$} The cohomology of $X$ and of $X_{\mathrm{red}}$ agree, so let $X$ be reduced. Since $\Fq$ is perfect,
$X$ has a dense open subscheme that is smooth, and on a dense open subscheme of it $\F$ is lisse; shrinking, there is a
dense open $U\subset X$, smooth and affine of pure dimension on each connected component, with $\F|_U$ lisse. Its
complement $Z$ has smaller dimension. By~\ref{cit:local}, every eigenvalue of $F$ on $H^i_c(X,\F)$ is an eigenvalue on
$H^i_c(U,\F)$ or on $H^i_c(Z,\F)$. The first are bounded by Step~1 (applied to each component of $U$), the second by
induction on $\dim X$, since $\F|_Z$ again has pointwise weight $\le w$. Dimension $0$ is covered by Step~1.
\end{proof}
\uses{Theorem~\ref{thm:middle}; Lemma~\ref{lem:good}; \ref{cit:artin}, \ref{cit:duality}, \ref{cit:local}.}

\begin{corollary}[$=$ Corollary~\ref{cor:B}]\label{cor:pure}
Let $X$ be smooth and proper over $\Fq$ of pure dimension $d$, and $\F$ lisse and pointwise pure of weight $w$. Then every
eigenvalue $\alpha$ of $F$ on $H^i(X,\F)$ satisfies $|\iota\alpha|=q^{(w+i)/2}$.
\end{corollary}
\begin{proof}
Since $X$ is proper, $H^i(X,\F)=H^i_c(X,\F)$, so $|\iota\alpha|\le q^{(w+i)/2}$ by Theorem~\ref{thm:all}. The dual
$\F^\vee$ is lisse and pointwise pure of weight $-w$; this uses purity at every point, not only the upper bound, since the
eigenvalues of $\Frob_x$ on $\F^\vee_{\bar x}$ are the inverses of those on $\F_{\bar x}$. By~\ref{cit:duality} the pairing
$H^i(X,\F)\otimes H^{2d-i}(X,\F^\vee)\to\Ql(-d)$ is perfect and Frobenius-equivariant, so $q^d/\alpha$ is an eigenvalue of
$F$ on $H^{2d-i}(X,\F^\vee)$. Theorem~\ref{thm:all} for $\F^\vee$ gives $|\iota(q^d/\alpha)|\le q^{(-w+2d-i)/2}$, that is,
$|\iota\alpha|\ge q^{(w+i)/2}$. For $\F=\Ql$ the hypothesis holds with $w=0$ for every $\iota$.
\end{proof}
\uses{Theorem~\ref{thm:all}; \ref{cit:duality}.}

\begin{remark}[Over a base]
For $f:X\to S$ separated of finite type over $\Fq$ and $s$ a closed point of $S$, the stalk of $R^if_!\F$ at $\bar s$ is
$H^i_c(X_{\bar s},\F)$ by proper base change \cite[XVII, \S5]{SGA4}. So Theorem~\ref{thm:all}, applied to the fibre
$X_s$ over $\kappa(s)$, bounds the Frobenius eigenvalues of $R^if_!\F$ at $s$ by $|\kappa(s)|^{(w+i)/2}$.
\end{remark}

\begin{remark}[The sign, and Grothendieck's evanescent part]
The ``definite sign'' of Alp\"oge's remark \cite{Alpoge} is the parity $(-1)^{D-1}$ of Proposition~\ref{prop:remainder}:
after this sign the remainder of the point count is a power sum with positive multiplicities. The module $K$ is the
compactly supported, affine form of the part of the cohomology of a hyperplane section that the Gysin map kills. In
Grothendieck's unpublished notes in the Montpellier archive (cote 16, scans 034 and 036 in the transcription of
\cite{Archive}), this part, $\operatorname{Ker}(H^{n-1}(Y)\to H^{n+1}(X))$, is his \emph{partie \'evanescente}
$E^{n-1}$. The identification is ours and plays no role in the proof.
\end{remark}

% ==================================================================================
\fixedhead{Limits and open points}
% ==================================================================================

\begin{itemize}[leftmargin=*]
\item The Lefschetz input is~\ref{cit:katz}, Katz's Appendix A.5, whose proof was read at the source by the independent
checker and found to use no weights; it was also re-proved independently (Verification record). The printed proof of
Katz's Corollary 3.4.1 goes through the Fourier transform (Brylinski, Katz--Laumon) and is not used here.
\item The other cited results (Section~\ref{sec:inputs}) are standard \'etale cohomology from SGA 4 and SGA $4\frac12$;
their locators are those used by the producing run, and they were not reopened for this text.
\item Nothing is formalized.
\end{itemize}

% ==================================================================================
\fixedhead{Verification record}
% ==================================================================================

\noindent
\begin{tabularx}{\textwidth}{@{}L{3.4cm}L{2.0cm}Y@{}}
\toprule
Who & Date & Type, scope and verdict\\
\midrule
GPT-6 Astra (ChatGPT Pro mode), the producing run & 21 Nov 2026 & In-run checks by the run's sub-agents (209 check
events in its CHECK\_EVENTS.tsv). Author-side.\\
Claude Opus 5.5 at effort high (intake) & 21 Nov 2026 & Targeted proof check of every step of Lemmas 3.1--3.3,
Proposition 4.1, Theorem 5.1, Theorem 6.1 and Corollary 6.2 (the run's HS-01--HS-03, SG-08--SG-11); Katz's Corollary
3.4.1 and Appendix A.1--A.3 read at the source. Confirmed, no gap.\\
Claude Opus 5.5 at effort max (main session) & 21 Nov 2026 & Rederived Lemma 3.3, the identity of Proposition 4.1 with its
twists and signs, and the exponent in Theorem 5.1, Step 3; read Katz's Theorem 3.4, Corollaries 3.4.1 and 3.4.3 and
Lemma 3.5 at the source. No fault found.\\
Claude Fable 5.2 at effort max (independent check; not shown the other checks) & 21 Nov 2026 & Lemmas 3.1--3.3,
Proposition 4.1, Theorem 6.1 and Corollary 6.2 re-proved; Proposition 4.1 also checked on a worked example. Theorem 5.1
confirmed with one correction of the citation route: take the Lefschetz input from Katz's Appendix A.3--A.5 through the
graph immersion, not from the printed proof of Corollary 3.4.1. This text adopts the correction as Lemma 2.8, and the
checker re-proved that input itself. Bertini by a dimension count, also adopted. Katz pp.~152--158, 179--181 and 219--222
read at the source. No gap, no circular input. Not opened: SGA 4, SGA $4\frac12$, BBD, Brylinski, Katz--Laumon,
Jouanolou.\\
\bottomrule
\end{tabularx}

% ==================================================================================
\fixedhead{Provenance}
% ==================================================================================

The question was prompt 32 (HODGE\_SIGN) of our working-through of the two proofs of generalized Riemann hypotheses of
20 November 2026, written by Claude Opus 5.5 at effort max on 20 Nov 2026. The proof is Part I, Section 1 of the run's
return (GPT-6 Astra, ChatGPT Pro mode, 21 Nov 2026, 00:14--00:56 UTC). The prompt as sent, the full return, the intake and
the check records are public at
\url{https://github.com/KokunoYumeto/grothendieck-montpellier-archive/tree/main/research/grh_20261120/deligne_sign_only}
and \url{https://github.com/KokunoYumeto/grothendieck-montpellier-archive/tree/main/raw/research_returns/HODGE_SIGN_final_20261121T0057Z_unzipped}.
This text was written from the run's paper by Claude Opus 5.5 at effort max; the statements and proofs follow the run's,
restated and with the steps written out.

% ==================================================================================
\fixedhead{Declarations}
% ==================================================================================

\noindent\textbf{AI use.} The proof was produced by GPT-6 Astra (ChatGPT Pro mode); the text was written by Claude (Claude
Opus 5.5 at effort max); the checks are listed in the verification record. No part of the text was written by a human.
\textbf{Human understanding.} None. \textbf{Licence.} CC0 1.0. \textbf{Credits.} The idea of a sign-only proof from
$qN=M^k+R$ is L.~Alp\"oge's (footnote 3 of \cite{Alpoge}); the generic weak Lefschetz theorem is N.~Katz's
\cite{KatzACT}, and the proof of it in his appendix is P.~Deligne's. \textbf{Version.} 1.

\begin{thebibliography}{99}
\bibitem{Alpoge} L.~Alp\"oge, note on the Riemann hypothesis, 20 November 2026, \url{https://alpo.ge/remark.pdf}.
Footnote 3. Opened.
\bibitem{WeilI} P.~Deligne, La conjecture de Weil.~I, Publ. Math. IH\'ES 43 (1974), 273--307. Th\'eor\`eme 1.6 and
(1.5.4). Not reopened.
\bibitem{WeilII} P.~Deligne, La conjecture de Weil.~II, Publ. Math. IH\'ES 52 (1980), 137--252. Th\'eor\`eme 3.3.1. Not
reopened.
\bibitem{KatzACT} N.~M.~Katz, Affine cohomological transforms, perversity, and monodromy, J. Amer. Math. Soc. 6 (1993),
149--222. Theorem 3.4, Corollaries 3.4.1 and 3.4.3, Lemma 3.5, Appendix A.3--A.5, pp.~152, 180--182, 219--221. Opened
(\url{https://web.math.princeton.edu/~nmk/old/ACT.pdf}).
\bibitem{Laumon} G.~Laumon, Transformation de Fourier, constantes d'\'equations fonctionnelles et conjecture de Weil, Publ.
Math. IH\'ES 65 (1987), 131--210. Not opened.
\bibitem{Kedlaya} K.~S.~Kedlaya, Fourier transforms and $p$-adic ``Weil II'', Compos. Math. 142 (2006), 1426--1450. Not
opened.
\bibitem{Scholl} A.~J.~Scholl, Hypersurfaces and the Weil conjectures, \href{https://arxiv.org/abs/0912.2073}{arXiv:0912.2073}.
Introduction. Opened.
\bibitem{Jouanolou} J.-P.~Jouanolou, Th\'eor\`emes de Bertini et applications, Progr. Math. 42, Birkh\"auser, 1983.
Th\'eor\`eme I.6.10. Not opened.
\bibitem{SGA4} M.~Artin, A.~Grothendieck, J.-L.~Verdier, Th\'eorie des topos et cohomologie \'etale des sch\'emas (SGA 4),
Lecture Notes in Math. 269, 270, 305, Springer, 1972--1973. Expos\'es XIV, XVII, XVIII. Not reopened.
\bibitem{SGA4h} P.~Deligne, Cohomologie \'etale (SGA $4\frac12$), Lecture Notes in Math. 569, Springer, 1977. Rapport,
Th\'eor\`eme 3.2. Not reopened.
\bibitem{Stacks} The Stacks Project, Tag 00OY, \url{https://stacks.math.columbia.edu/tag/00OY}. Not reopened.
\bibitem{Archive} Our transcription of Grothendieck's Montpellier archive,
\url{https://github.com/KokunoYumeto/grothendieck-montpellier-archive}. Cote 16. Opened.
\end{thebibliography}

\end{document}
